\documentclass[10pt,a4paper]{amsart}
\usepackage{amsmath,amssymb,amsthm,graphicx}
% Retained minimal preamble for the Stacks Project chapter coherent.tex (Stacks chapter, pre-delegation retained fragment for stacks-coherent).
% Verbatim excerpt of the chapter preamble (sources/stacks-spaces-descent/preamble.tex): the reference/slogan/history
% environment block (native lines 7-12), Xy-pic setup (native lines 14-19) and the native theorem declarations with
% the mathematical macros (native lines 154-217). Copied mechanically; nothing added, removed or redefined. The
% wrapper class and the amsmath/amssymb/amsthm/graphicx base packages are supplied by the assembler, so the
% class/xr-hyper/hyperref/externaldocument part of the native preamble is omitted.
% For dealing with references we use the comment environment
\usepackage{verbatim}
\newenvironment{reference}{\comment}{\endcomment}
%\newenvironment{reference}{}{}
\newenvironment{slogan}{\comment}{\endcomment}
\newenvironment{history}{\comment}{\endcomment}
% For commutative diagrams we use Xy-pic
\usepackage[all]{xy}

% We use 2cell for 2-commutative diagrams.
\xyoption{2cell}
\UseAllTwocells
% Theorem environments.
%
\theoremstyle{plain}
\newtheorem{theorem}[subsection]{Theorem}
\newtheorem{proposition}[subsection]{Proposition}
\newtheorem{lemma}[subsection]{Lemma}

\theoremstyle{definition}
\newtheorem{definition}[subsection]{Definition}
\newtheorem{example}[subsection]{Example}
\newtheorem{exercise}[subsection]{Exercise}
\newtheorem{situation}[subsection]{Situation}

\theoremstyle{remark}
\newtheorem{remark}[subsection]{Remark}
\newtheorem{remarks}[subsection]{Remarks}

\numberwithin{equation}{subsection}

% Macros
%
\def\lim{\mathop{\mathrm{lim}}\nolimits}
\def\colim{\mathop{\mathrm{colim}}\nolimits}
\def\Spec{\mathop{\mathrm{Spec}}}
\def\Hom{\mathop{\mathrm{Hom}}\nolimits}
\def\Ext{\mathop{\mathrm{Ext}}\nolimits}
\def\SheafHom{\mathop{\mathcal{H}\!\mathit{om}}\nolimits}
\def\SheafExt{\mathop{\mathcal{E}\!\mathit{xt}}\nolimits}
\def\Sch{\mathit{Sch}}
\def\Mor{\mathop{\mathrm{Mor}}\nolimits}
\def\Ob{\mathop{\mathrm{Ob}}\nolimits}
\def\Sh{\mathop{\mathit{Sh}}\nolimits}
\def\NL{\mathop{N\!L}\nolimits}
\def\CH{\mathop{\mathrm{CH}}\nolimits}
\def\proetale{{pro\text{-}\acute{e}tale}}
\def\etale{{\acute{e}tale}}
\def\QCoh{\mathit{QCoh}}
\def\Ker{\mathop{\mathrm{Ker}}}
\def\Im{\mathop{\mathrm{Im}}}
\def\Coker{\mathop{\mathrm{Coker}}}
\def\Coim{\mathop{\mathrm{Coim}}}

% Boxtimes
%
\DeclareMathSymbol{\boxtimes}{\mathbin}{AMSa}{"02}

%
% Macros for moduli stacks/spaces
%
\def\QCohstack{\mathcal{QC}\!\mathit{oh}}
\def\Cohstack{\mathcal{C}\!\mathit{oh}}
\def\Spacesstack{\mathcal{S}\!\mathit{paces}}
\def\Quotfunctor{\mathrm{Quot}}
\def\Hilbfunctor{\mathrm{Hilb}}
\def\Curvesstack{\mathcal{C}\!\mathit{urves}}
\def\Polarizedstack{\mathcal{P}\!\mathit{olarized}}
\def\Complexesstack{\mathcal{C}\!\mathit{omplexes}}
% \Pic is the operator that assigns to X its picard group, usage \Pic(X)
% \Picardstack_{X/B} denotes the Picard stack of X over B
% \Picardfunctor_{X/B} denotes the Picard functor of X over B
\def\Pic{\mathop{\mathrm{Pic}}\nolimits}
\def\Picardstack{\mathcal{P}\!\mathit{ic}}
\def\Picardfunctor{\mathrm{Pic}}
\def\Deformationcategory{\mathcal{D}\!\mathit{ef}}
\usepackage{cleveref}
\crefname{originallemma}{original lemma}{original lemmas}
\Crefname{originallemma}{Original lemma}{Original lemmas}
\makeatletter
\newlabel{algebra-lemma-Noetherian-power}{{00IM}{source}{}{sourcepointer-1}{}}
\newlabel{algebra-lemma-Noetherian-power@cref}{{[originallemma][0][]00IM}{[1][0][]source}{}{}{}}
\makeatother
\crefname{originallemma}{original lemma}{original lemmas}
\Crefname{originallemma}{Original lemma}{Original lemmas}
\makeatletter
\newlabel{algebra-lemma-change-ideal-completion}{{0319}{source}{}{sourcepointer-2}{}}
\newlabel{algebra-lemma-change-ideal-completion@cref}{{[originallemma][0][]0319}{[1][0][]source}{}{}{}}
\makeatother
\crefname{originallemma}{original lemma}{original lemmas}
\Crefname{originallemma}{Original lemma}{Original lemmas}
\makeatletter
\newlabel{algebra-section-spectrum-ring}{{00DY}{source}{}{sourcepointer-3}{}}
\newlabel{algebra-section-spectrum-ring@cref}{{[originallemma][0][]00DY}{[1][0][]source}{}{}{}}
\makeatother
\crefname{originallemma}{original lemma}{original lemmas}
\Crefname{originallemma}{Original lemma}{Original lemmas}
\makeatletter
\newlabel{lemma-inverse-systems-affine}{{087W}{source}{}{sourcepointer-4}{}}
\newlabel{lemma-inverse-systems-affine@cref}{{[originallemma][0][]087W}{[1][0][]source}{}{}{}}
\makeatother
\title{Formal systems depend only on the support: for X Noetherian and I, J quasi-coherent sheaves of ideals with the same vanishing closed subset, the categories Coh(X, I) and Coh(X, J) are equivalent.}
\author{The Stacks Project Authors}
\date{stacks-project:coherent.tex:local-snapshot-20261012; GNU FDL 1.2 or later}
\begin{document}
\maketitle
\noindent\textit{Editorial note.} Statement span is lines 6922-6930 of sources/stacks-coherent/pinned-original-source.tex; proof span is lines 6932-6960, the complete author proof (27 body lines).\par\medskip
\noindent\textit{Editorial note.} Navigation stubs: 4 cited-but-not-included labels are supplied as non-load-bearing pointers with their pinned tags where the snapshot tag index has them.\par\medskip
\noindent\textit{Editorial note.} extra\_packages is empty on purpose: the retained preamble already loads xy with options and a second declaration clashes in the assembler wrapper.\par\medskip
\noindent\textit{Editorial note.} Text is retained verbatim, including the author's own typos and placeholders.\par\medskip
\noindent\textit{Editorial note.} Pool stacks-coherent backfill wave 33d. precard\_gate RESULT=OK (probe L9999).\par\medskip
\noindent\textit{Allowed standard prerequisites (editorial).}\begin{itemize}
\item Explicit prerequisite boundary. See why-field; every pointer named here occurs verbatim in the retained statement/proof text.
\end{itemize}\medskip
\ifdefined\hypertarget\hypertarget{sourcepointer-1}{}\fi
\noindent\textit{Source reference.} Original lemma 00IM: sources/stacks-coherent/pinned-original-source.tex, source lines 6960--6960. Cited by the retained statement or proof but not part of the collected spans; supplied only as a navigation pointer. The argument does not depend on its content.\par\medskip
\ifdefined\hypertarget\hypertarget{sourcepointer-2}{}\fi
\noindent\textit{Source reference.} Original lemma 0319: sources/stacks-coherent/pinned-original-source.tex, source lines 6960--6960. Cited by the retained statement or proof but not part of the collected spans; supplied only as a navigation pointer. The argument does not depend on its content.\par\medskip
\ifdefined\hypertarget\hypertarget{sourcepointer-3}{}\fi
\noindent\textit{Source reference.} Original lemma 00DY: sources/stacks-coherent/pinned-original-source.tex, source lines 6960--6960. Cited by the retained statement or proof but not part of the collected spans; supplied only as a navigation pointer. The argument does not depend on its content.\par\medskip
\ifdefined\hypertarget\hypertarget{sourcepointer-4}{}\fi
\noindent\textit{Source reference.} Original lemma 087W: sources/stacks-coherent/pinned-original-source.tex, source lines 6489--6495. Cited by the retained statement or proof but not part of the collected spans; supplied only as a navigation pointer. The argument does not depend on its content.\par\medskip
\begin{lemma}
\label{lemma-inverse-systems-ideals-equivalence}
Let $X$ be a Noetherian scheme. Let
$\mathcal{I}, \mathcal{J} \subset \mathcal{O}_X$
be quasi-coherent sheaves of ideals.
If $V(\mathcal{I}) = V(\mathcal{J})$ is the same closed subset
of $X$, then $\textit{Coh}(X, \mathcal{I})$ and $\textit{Coh}(X, \mathcal{J})$
are equivalent.
\end{lemma}

\begin{proof}
First, assume $X = \Spec(A)$ is affine. Let $I, J \subset A$ be the ideals
corresponding to $\mathcal{I}, \mathcal{J}$. Then $V(I) = V(J)$
implies we have $I^c \subset J$ and $J^d \subset I$ for some $c, d \geq 1$
by elementary properties of the Zariski topology
(see Algebra, Section \ref{algebra-section-spectrum-ring} and
Lemma \ref{algebra-lemma-Noetherian-power}).
Hence the $I$-adic and $J$-adic completions of $A$ agree, see
Algebra, Lemma \ref{algebra-lemma-change-ideal-completion}.
Thus the equivalence follows from Lemma \ref{lemma-inverse-systems-affine}
in this case.

\medskip\noindent
In general, using what we said above and the fact that
$X$ is quasi-compact, to choose $c, d \geq 1$ such that
$\mathcal{I}^c \subset \mathcal{J}$ and $\mathcal{J}^d \subset \mathcal{I}$.
Then given an object $(\mathcal{F}_n)$ in
$\textit{Coh}(X, \mathcal{I})$ we claim that the
inverse system
$$
(\mathcal{F}_{cn}/\mathcal{J}^n\mathcal{F}_{cn})
$$
is in $\textit{Coh}(X, \mathcal{J})$. This may be checked on
the members of an affine covering; we omit the details.
In the same manner we can construct an object of
$\textit{Coh}(X, \mathcal{I})$ starting with an object of
$\textit{Coh}(X, \mathcal{J})$. We omit the verification
that these constructions define mutually quasi-inverse functors.
\end{proof}

\end{document}
